\section{Verification, remaining estimates, and interpretation}\label{sec:verification}

\subsection{Three different forms of evidence}
A written proof, an accepting finite computation, and a byte-identical publication answer different questions. A proof must justify its infinite sums, analytic continuations, and uniformity. A finite checker can establish algebraic identities, a directed finite inequality, or an existence certificate on an explicitly covered parameter cube. A manifest or Git object identity establishes which bytes were tested or published. None substitutes for either of the others.

The theta certificates in Section~\ref{sec:theta} are genuine finite existence claims in the source: their accepting argument encloses the full analytic moment source and the full finite parameter cube. By contrast, many arithmetic checkers test only finite algebra attached to an infinite proof. The numbers of predicates in their reports include repeated domain and setup guards. They are not numbers of independent theorems, prime ranges, or zeta zeros verified.

During preparation of this manuscript, the original local packets for \cite{R1000,R1002,R1004} were replayed without editing their scientific files. The complete positive-determinant replay reproduced its ordinary and optimized 3,370-predicate reports and all eight designated corrupted executions. The five-route replay authenticated fourteen manifested payloads, reproduced 20,279 predicates per mode, and rejected twelve designated corruptions. Both completed with empty standard error. The failed-$0.87$ audit reproduced its archived report byte for byte in both isolated ordinary and optimized Python. Its report is an exponent-model calculation, not a zeta-zero computation.

A separate small manuscript checker recomputes the 850 rational middle-class comparisons in \eqref{eq:271-margin}, the scalar constants of the failed comparison, the displayed residual-length ranges, the conditional recurrence exponents, the native frequency cutoff, the graph barrier, and the inequalities using the rounded bounds in Table~\ref{tab:theta-certificates}. Rechecking those rounded bounds does not reconstruct the theta source or prove its enclosure. The full \cite{R1003} certificate campaigns, the \cite{R989} analytic contracts, and the upstream reflection and local-domain arguments in \cite{R977} were \emph{not} newly executed or independently certified by compiling this manuscript. Their exact pins and reading scopes are part of the provenance ledger.

Some source packets record separate AI-agent reviews within their originating session. Those reports remain source records; they are not external human referee reports, and they are not rebranded here as new independent reviews. This manuscript was synthesized with substantial AI assistance, including source reading, mathematical exposition, code, and typesetting. No independent mathematical referee, Lean proof, repository-wide validation, actual $L$-value census, or zeta-zero computation is claimed for its preparation. Human author approval, mathematical review, and arXiv submission remain separate actions.

The distributed LaTeX source uses a conventional \texttt{amsart} build, fixed dates, local section files, and a pre-generated bibliography. It requires no shell escape or online data access to compile. The proof-source ledger and build instructions accompany the repository version. The arXiv source archive contains the necessary source files rather than the resulting PDF or transient build products. A manuscript compile and visual inspection verify presentation, not mathematical correctness.

\subsection{The estimates not supplied}
\paragraph{The ordinary scalar residual.}
Equation~\eqref{eq:remaining-scalar} removes the auxiliary field variance without removing $\abs{\zeta(s)M_Z(s)-1}^2$. A strict upper bound below one at every point in a desired region would exclude ordinary zeros, but the source family approximation does not provide it. Similarly, the signed field mean in Theorem~\ref{thm:field-average} has not been bounded at the power required by its zero-free consumer.

\paragraph{The native arithmetic contraction.}
Equations~\eqref{eq:rough-comparison} and \eqref{eq:frequency-tail-general} pay specific complete sectors. The remaining large-core or physically restricted low-frequency amplitude still needs an arithmetic estimate with a genuinely subquadratic power of the old energy, or another estimate of comparable strength. The synthetic counterexample shows why generic coefficient bounds and balance identities alone cannot supply the proposed contraction.

\paragraph{Unbounded theta feasibility.}
The finite certificates establish matching through order 26. Theorem~\ref{thm:weighted-canonical} would use positive weighted solutions at an unbounded set of orders. Neither two successful root-extension steps nor finitely many positive Hankel matrices establishes that premise. In the strip programme, controlled fields and negative edges enlarge the admissible class but have not produced a cofinal theta-matching sequence.

\paragraph{A signed filtered or determinant estimate.}
The two incommensurate filters preserve every nonreal target pole, but no improved full physical remainder estimate is proved. The all-mass determinant theorem reduces conditioning and pays the complete long heat tail, but not the centered short-time trace or height. These are different representations of a remaining arithmetic task, not automatic deductions from source identification.

\paragraph{Independent review of the broad-field proposal.}
The displayed $271/300$ composition retains the scope and dependencies of \cite{R977}. Its review requires reconstructing the short-support estimate, ramified local correction, selected-character interfaces, and all contour tails. Finite exponent margins, including the freshly recomputed \eqref{eq:271-margin}, cannot validate those analytic premises.

\subsection{Conclusion}
The recent contributions are not exhausted by a sequence of decimals near one. They include a uniform actual-field coefficient comparison; a full-$L$ approximation and an exact auxiliary-variance optimization; a two-polynomial transfer for capped reciprocal laws; complete denominator, core, and frequency-sector bounds; finite real-zero approximants through order 26; a less restrictive weighted canonical criterion; a controlled negative-coupling strip class; and explicitly regularized positive determinants. The failed-$0.87$ comparison is also part of the result record because it identifies a precise missing power saving in a specified estimate system.

What these contributions share is a requirement of source fidelity: common factors, square aliases, ramified zeros, prime powers, complete cross terms, physical cutoffs, and limit quantifiers are retained rather than discarded when inconvenient. They provide component theorems and more explicit targets. They do not yet provide the all-scale signed estimate or unbounded positive construction needed for a new global ordinary-zeta zero-free region or RH.
